\documentclass[11pt]{article}
\usepackage[margin=1in]{geometry}
\usepackage{amsmath,amssymb,amsthm,lmodern,microtype}
\usepackage[T1]{fontenc}
\usepackage{xurl}
\usepackage[colorlinks=true,urlcolor=blue,citecolor=blue,linkcolor=blue]{hyperref}
\newtheorem{theorem}{Theorem}
\newtheorem{lemma}{Lemma}
\newcommand{\C}{\mathbb C}
\newcommand{\dd}{\,dA}
\title{A fixed-area exceptional set for positive squared-pole sums}
\author{Alec Kriebel\\ \small Independent researcher\\
\small ORCID: \href{https://orcid.org/0009-0001-9320-500X}{0009-0001-9320-500X}}
\date{17 October 2026 \quad Version 1.0.0}
\begin{document}
\maketitle
\begin{abstract}
For every finite positive sum \(g(z)=\sum_j\lambda_j(z-z_j)^{-2}\)
with \(\sum_j\lambda_j=1\), we construct a compact set \(S(g)\) of
area exactly \(\pi\) such that
\[
 \int_{\{|z|<R\}\setminus S(g)}|g(z)|\dd
 \le 2\pi\log R+\pi(\log2+e^{-1})\qquad(R>1).
\]
The same set serves every radius, with no restriction on the poles.
This gives an affirmative deduction for Baernstein's question,
Hayman--Lingham Problem~7.37, from the established
Eremenko--Hamilton indicator estimate and Levine--Peres quadrature-set
construction. The logarithmic coefficient is sharp; no optimality is
claimed for the additive constant. Historical first priority is not
certified. AI tools were used extensively; this note is unrefereed.
\end{abstract}

\section{The question and the result}
Write \(D_r=\{z\in\C:|z|<r\}\) and use unnormalized planar Lebesgue
area \(dA=dx\,dy\). Baernstein's question in \cite[Problem~7.37,
printed pp.~171--172]{HL} asks for an absolute constant \(C\) and,
for each positive normalized squared-pole sum, an area-\(\pi\) set
independent of \(R\), with integral at most \(2\pi\log R+C\).
Here is an explicit answer.
\begin{theorem}\label{thm}
Let \(n\ge1\), \(z_j\in\C\), \(\lambda_j>0\), and
\(\sum_{j=1}^n\lambda_j=1\), and set
\(g(z)=\sum_{j=1}^n\lambda_j/(z-z_j)^2\).
There exists a compact \(S=S(g)\), of area exactly \(\pi\), such that,
simultaneously for every real \(R>1\),
\begin{equation}\label{main}
 \int_{D_R\setminus S}|g|\dd
 \le 2\pi\log R+\pi(\log2+e^{-1}).
\end{equation}
These are ordinary finite nonnegative Lebesgue integrals. Poles may
coincide. No smaller logarithmic coefficient works with a finite
absolute additive constant.
\end{theorem}

\section{The established inputs}
Use the normalization
\[
 Tf(z)=-\frac1\pi\operatorname{pv}
       \int_{\C}\frac{f(w)}{(w-z)^2}\,dA(w).
\]
For bounded-set indicators this has its usual \(L^2\) realization.
Off the closed support the integral converges absolutely.
Eremenko--Hamilton \cite[printed p.~2794, proof pp.~2796--2797]{EH}
prove, for measurable \(E\subset D_1\),
\begin{equation}\label{EH}
 \int_{D_1\setminus E}|T1_E|\dd
 \le |E|\log\frac{\pi}{|E|}.
\end{equation}
The zero-area value is defined by continuity. Their compact
capacity-one ambient-set formulation includes the closed unit disk;
disk boundaries are null. Measurable indicators follow also from
compact approximation, \(L^2\) continuity and Fatou.
Scaling \(z=L\zeta\) gives
\begin{equation}\label{scale}
 \int_{D_L\setminus E}|T1_E|\dd
 \le a\log\frac{\pi L^2}{a},\qquad
 E\subset D_L,\quad a=|E|.
\end{equation}
Indeed \(T1_E(L\zeta)=T1_{E/L}(\zeta)\).
The indicator estimate uses the conformal-on-\(E\) case of
\cite[Theorem~1(i)]{EH}. The later correction to its combined
general-distortion statement does not change this case
\cite[Theorem~5 and correction discussion]{EH2002}.

The second input is the following consequence of the smash-sum
construction in Levine--Peres \cite{LP}. Given finitely many
points \(a_j\) and positive masses \(m_j\), there is a bounded open
\(\Omega\), containing each disk \(B_j\) of area \(m_j\) centered at
\(a_j\), such that
\begin{equation}\label{quad}
 |\partial\Omega|=0,\qquad |\Omega|=\sum_jm_j,\qquad
 \int_\Omega h\dd=\sum_jm_jh(a_j)
\end{equation}
for every real integrable harmonic function on the whole of \(\Omega\).
Connectedness is unnecessary. Precise locators in the inspected
1 May 2009 author manuscript are definition~(4), Proposition~2.12,
Corollary~2.13, Lemma~2.15, Lemma~6.1 and Proposition~6.6
(pp.~6, 12--17, 54--55, 58--59).
Proposition~6.6 gives the superharmonic inequality; apply it to
\(h\) and \(-h\).

Here the regularity hypotheses allow finite iteration:
at each binary step the sum of two domain indicators is bounded,
compactly supported, continuous almost everywhere off null boundaries,
and takes values \(0,1,2\). The density-gap condition of
Proposition~2.12 holds with any \(0<\lambda<1\). That proposition
gives a null boundary, Corollary~2.13 gives area additivity, and
Lemma~2.15 gives boundedness. Containment, associativity and
Proposition~6.6 give \eqref{quad} for the finite smash sum.
The disk mean-value identity supplies the point evaluations.
These are established quadrature and partial-balayage results,
with antecedents credited in \cite{LP}; their existence or regularity
is not claimed as new here.

\section{One set and its entire exterior field}
Combine repeated poles by adding their weights. Apply \eqref{quad}
with \(a_j=z_j\) and \(m_j=\pi\lambda_j\), and put
\(S=\overline\Omega\). Then \(S\) is compact and
\[
 |S|=|\Omega|=\pi.
\]
It contains an open neighborhood of every pole. Fix \(z\notin S\).
The real and imaginary parts of \(w\mapsto(w-z)^{-2}\) are bounded
integrable harmonic functions on \(\Omega\), so \eqref{quad} gives
\begin{equation}\label{field}
 \int_\Omega\frac{dA(w)}{(w-z)^2}
  =\pi\sum_j\frac{\lambda_j}{(z_j-z)^2}
  =\pi g(z),\qquad T1_\Omega(z)=-g(z).
\end{equation}
The identity holds on every component of \(\C\setminus S\), including
any bounded holes. Harmonic tests are on the whole domain, so no
analytic continuation from the unbounded component is being assumed.
All these integrals are absolutely convergent.
Moreover, outside \(S\), \(|z-z_j|\ge\sqrt{\lambda_j}\), hence
\(|g(z)|\le n\). This already proves local finiteness.
The set \(S\) is fixed before any \(R\) is chosen.

\section{Localization and the uniform constant}
Fix \(R>1\). Set
\[
 L=\sqrt2R,\qquad E=\Omega\cap D_L,\qquad
 F=\Omega\setminus D_L,\qquad a=|E|\in[0,\pi].
\]
By \eqref{field}, linearity and \eqref{scale},
\begin{equation}\label{near}
 \int_{D_R\setminus S}|g|\dd
 \le a\log\frac{2\pi R^2}{a}
       +\int_{D_R\setminus S}|T1_F|\dd,
\end{equation}
since \(D_R\setminus S\subset D_L\setminus E\).
For \(|w|=\rho>R\), polar integration gives
\begin{equation}\label{kernel}
 \int_{D_R}\frac{dA(z)}{|z-w|^2}
 =\int_0^R\frac{2\pi r}{\rho^2-r^2}\,dr
 =\pi\log\frac{\rho^2}{\rho^2-R^2}.
\end{equation}
The angular identity follows by the geometric series for the Poisson
kernel. For \(w\in F\), \(\rho\ge\sqrt2R\), so \eqref{kernel}
is at most \(\pi\log2\). The positive separation \(L-R\) justifies
the ordinary transform integral throughout \(D_R\). Tonelli yields
\[
 \int_{D_R\setminus S}|T1_F|\dd
 \le\frac1\pi\int_F\int_{D_R}\frac{dA(z)\,dA(w)}{|z-w|^2}
 \le(\pi-a)\log2.
\]
Consequently,
\begin{align*}
 \int_{D_R\setminus S}|g|\dd
 &\le a\log\frac{2\pi R^2}{a}+(\pi-a)\log2\\
 &=2a\log R+\pi\log2+a\log\frac\pi a\\
 &\le2\pi\log R+\pi\log2+\frac\pi e .
\end{align*}
The last inequality uses \(\log R>0\), \(a\le\pi\) and
\(\max_{0\le t\le1}t\log(1/t)=1/e\).
Continuous endpoint values cover \(a=0\) and \(a=\pi\).
Only the auxiliary split depends on \(R\). This proves
\eqref{main} for the same \(S\) at every radius.

For sharpness of the leading coefficient, take \(g(z)=z^{-2}\).
For any measurable \(S\) of area \(\pi\), write
\(A=D_R\setminus D_1\). The integrand is at least \(1\) on \(D_1\)
and at most \(1\) on \(A\), while
\[
 |S\cap A|\le |S\setminus D_1|
   =|D_1\setminus S|.
\]
Thus the integral gained on \(D_1\setminus S\) dominates the finite
integral lost on \(S\cap A\), and
\[
 \int_{D_R\setminus S}|z|^{-2}\dd
 \ge \int_A |z|^{-2}\dd=2\pi\log R.
\]
This also holds when the left side is infinite. Letting \(R\to\infty\)
excludes a smaller coefficient with any finite additive constant.
No optimality of the additive constant in \eqref{main} is asserted.

\section{Attribution, priority limits and verification}
This note supplies the displayed target-specific deduction from
established results. It claims no new Beurling estimate, quadrature
construction, or general balayage/operator mechanism.
Related recent work includes Lin's partial-balayage principles
\cite{Lin} and Guerra's real-input Beurling weak-type estimates
\cite{Guerra}; those inspected operator conclusions do not themselves
state the fixed-set, all-radius assertion proved here.
The generic ingredients used above were available by 2009.
A bounded priority audit identified no verified earlier full-target
answer; worldwide first discovery/publication is not certified.
In particular, Leung preliminary reports titled \emph{On a fixed
area problem}, AMS \emph{Abstracts}~11 (1990), p.~181, and
13 (1992), p.~440, were located through primary indexes, but their
mathematical texts were inaccessible. Their relevance remains
unresolved. Other unavailable full/final texts and exact inspection
scopes are listed in the supporting audit. Neither historical open
labels nor repository timestamps establish first priority.

The source archive preserves the original proof, source identities,
three distinct mathematical reviews, priority reviews, the correction
check and Leung supplement. Supplementary standard-library programs
test kernel normalization, quadrature fields (including holes),
scaling, scalar bounds and deliberate wrong alternatives.
Floating calculations are diagnostics, not validated quadrature or
proof certificates; the theorem rests on the universal argument above
and its explicitly stated established inputs.
AI tools were used extensively in proof development, computation,
adversarial checking, literature research and writing.
This paper and its supplements are unrefereed and have not undergone
external human peer review, journal acceptance or formal
proof-assistant certification.

\begin{thebibliography}{9}
\bibitem{HL} W.~K.~Hayman and E.~F.~Lingham,
\emph{Research Problems in Function Theory}, arXiv:1809.07200v2 (2018),
Problem~7.37, PDF pp.~172--173.
\url{https://arxiv.org/abs/1809.07200v2}.
\bibitem{EH} A.~Eremenko and D.~H.~Hamilton,
On the area distortion by quasiconformal mappings,
\emph{Proc. Amer. Math. Soc.}~123 (1995), 2793--2797.
\href{https://doi.org/10.1090/S0002-9939-1995-1283548-8}{doi:10.1090/S0002-9939-1995-1283548-8}.
\bibitem{LP} L.~Levine and Y.~Peres,
Scaling limits for internal aggregation models with multiple sources,
\emph{J. Anal. Math.}~111 (2010), 151--219.
\href{https://doi.org/10.1007/s11854-010-0015-2}{doi:10.1007/s11854-010-0015-2}.
Inspected author manuscript, 1 May 2009:
\url{https://lionellevine.github.io/scalinglimit.pdf};
versioned preprint \url{https://arxiv.org/abs/0712.3378v2}.
\bibitem{EH2002} A.~Eremenko and D.~H.~Hamilton,
Area distortion of quasiconformal mappings,
in \emph{Handbook of Complex Analysis: Geometric Function Theory},
vol.~1 (2002), 147, 149--160.
\href{https://doi.org/10.1016/S1874-5709(02)80007-5}{doi:10.1016/S1874-5709(02)80007-5}.
Relevant publisher preview inspected; full chapter unavailable.
\bibitem{Lin} Y.~Lin, \emph{Two partial balayage principles for
weak-type estimates}, author-hosted manuscript, inspected 17 October 2026.
\url{https://coolrmal.github.io/assets/papers/two-partial-balayage-principles.pdf}.
Body also pinned to Git version dated 5 October 2026 in the supplement;
historical public visibility is not certified.
\bibitem{Guerra} A.~Guerra,
\emph{Quasiconvexity of the Burkholder function on symmetric matrices},
arXiv:2608.23388v1 (24 August 2026).
\url{https://arxiv.org/abs/2608.23388v1}.
\end{thebibliography}
\end{document}
